\section{OUTDATED: The Hands-Off O2H Lemma}
  All formulations of the O2H lemma involve either aborting the distinguisher at a
    random point and measuring their query (which essentially amounts to
    guessing the query), or performing gentle measurements to check for bad
    events (which ever-so-slightly disturbs the combined state of the distinguisher and
    oracle for every query). Either approach leads to a loss that
    depends on the number of queries $q$.
    But what if we don't abort at all? Now that we have compressed oracles in
    our toolbox, can we instead just let $\distinguisher$ run until the
    end and then measure the database?
  
  The following is my attempt at an exact analysis of the best advantage
    an extractor, simulating a random oracle as a CRO for an optimal
    distinguisher, can ever get. As we will see, we will get very
    close to tightness, but then the inevitable
    roadblock will appear, leading us to the ubiquitous square root loss. 
    Still, the bound we achieve is indeed independent of $q$.
    Further you may agree that seeing what the nature of this roadblock is in
    itself quite informative,
    and it may even open a path towards \emph{proving} that the bound is optimal
    (at least for non-aborting reductions).

  %\begin{remark}
  %  Some differences to the Semi-Classical O2H lemma (SC-O2H): In this version, the 
  %    extractor $\extractor$ provides an \emph{independent}
  %    simulation of a random oracle, rather than one that is equal to the
  %    ``real'' oracle on all but
  %    $\ell$ points, and the adversary has to guess whether it
  %    gets the real, ``monolithic'' random oracle $\RO$, or the independently
  %    simulated random oracle $\RO[G]$, given some input
  %    $z$ which depends arbitrarily on $\ell$ points in the real random
  %    oracle $(\RO(x_1), \ldots, \RO(x_\ell))$. This matches a situation often
  %    found in game hops. 

  %  The goal of the extractor $\extractor$ will be to output a set of input values
  %    such that the set intersects with the set $\{x_1, \ldots, x_n\}$ on at
  %    least one value, after which point a reduction can either guess which one it
  %    needs to output, or use a PCA oracle or similar to find the value it wants.

  %  Finally, as mentioned, instead of measuring a random oracle query $i \in
  %    [q]$, the distinguisher is run until the end, before the extractor measures
  %    the oracle database and produces its output; hence the name ``Hands-Off
  %    O2H''.
  %\end{remark}

  \hhnote{ 
    Mention that using this lemma will sometimes mean explicitly bounding the
    runtime of $\adv$? (But surely this must be the case for normal O2H as
    well? Let first query be bad, and let $\adv$ enter an infinite loop upon
    noticing an inconsistency, making no further queries. If $i > 1$, then
    $\extractor$ never halts. Maybe this is often just taken to be implicit?
    Or $\adv$ is implicitly assumed to \emph{always} halt in polynomial time?
    Discuss when first introducing O2H.
  }

  \subsection{The Hands-Off O2H Lemma}
    \textbf{ERROR: The Hands-Off O2H Lemma is false, as shown by Joseph's
      counterexample, presented and analyzed in \autoref{sec:uncompute}.
      This section is kept around for historical reasons only,
      and as a jumping-off point for further investigations.}

    Before we turn to the main theorem, let us show how we can generically go from the
      situation considered therein, in which the $\RO$ and
      $\RO[G]$ are completely independent oracles, to something more akin to the
      Semi-Classical O2H, where the oracles only differ on input points in some set
      $\careset$. The reduction proving the lemma implements a strategy similar
      to (and inspired by) that of \autoref{lemma:two-equals-one}.

    \begin{lemma}
      Let $\careset \subseteq \bin^n$ be a set of input values $x^*_i$, and
        let $\RO, \RO[G]$ be (independent) random oracles from $\bin^n$ to $\bin^m$. Then
        there is a reduction $\bdv$ that, given $\careset$ and oracle access to
        $\RO$ and $\RO[G]$, provides a perfect
        simulation to any distinguisher $\ddv$ of a random oracle
        $\altRO$ satisfying
        \begin{enumerate}
          \item $\forall x \not\in \careset, \altRO(x) = \RO(x)$, 
          \item $\forall x \in \careset, \altRO(x) = \RO[G](x)$. 
        \end{enumerate}

      \noindent
      If $\adv$ makes $q$ queries to $\altRO$, then $\bdv$ makes $3q$ queries
        to $\RO$ and $2q$ queries to $\RO[G]$, or $5q$ queries in total.
        (If only independence from $\RO$ is required for the points in
        $\careset$, and not necessarily consistency
        with $\RO[G](x)$, then this can be reduced to $q$ queries to $\RO$ and
        $2q$ queries to $\RO[G]$, for $3q$ queries in total.)
    \end{lemma}

    \begin{proof}
      Let $\bdv$ keep an $m$-bit register $\register[Z]$, initialized to the
      all-zero state $\ket{0^m}_{\register[Z]}$. Let $\unitary_z$ be the
      unitary that acts on $\inputreg, \outputreg, \register[Z]$ such
      that, if $x \in \careset$, 
      \[ 
        \ket{x}\ket{y}\ket{z} \apply{\unitary_z} \ket{x}\ket{y \xor z}\ket{z}\, ,
      \]
      and otherwise acts like the identity.
      (This may be implemented using a series of controlled-controlled not gates
      (Toffoli gates), and ancilla qubits that store the results of equality
      check with each $x^*_i$ in $\careset$ individually, which may then be
      condenced to a single ancilla storing the result of the OR-operation on
      each of the individual ancillas.)

      Upon a query to $\altRO$ with
      registers $\inputreg, \outputreg$, $\bdv$ does the following:
      \begin{enumerate}
        \item Query $(\inputreg,\outputreg) \gets \RO(\inputreg, \outputreg)$,
        \item Query $(\inputreg,\register[Z]) \gets \RO(\inputreg, \register[Z])$,
        \item Query $(\inputreg,\register[Z]) \gets \RO[G](\inputreg, \register[Z])$, 
        \item $\applyU{\unitary_z}{\inputreg, \outputreg, \register[Z]}$,
        \item Query $(\inputreg,\register[Z]) \gets \RO[G](\inputreg, \register[Z])$,
        \item Query $(\inputreg,\register[Z]) \gets \RO(\inputreg, \register[Z])$,
        \item Return ${\inputreg, \outputreg}$ to $\ddv$.
      \end{enumerate}
      
      \noindent
      Let us see what the effect of the above operations are on the
        computational basis states $\ket x \ket y \ket 0^m$.

      \begin{align*}
        \ket{x}\ket{y}\ket{0^m} 
          \quad \apply{(1.)} \quad
          &\ket{x}\ket{y \xor \RO(x)}\ket{0^m}\, , \\
        \quad \apply{(2.)} \quad
          &\ket{x}\ket{y \xor \RO(x)}\ket{\RO(x)}\, , \\
        \quad \apply{(3.)} \quad
          &\ket{x}\ket{y \xor \RO(x)}\ket{\RO(x) \xor \RO[G](x)}\, .
        \intertext{
          If $x \neq \careset$, then the next steps give
        }
        \quad \apply{(4.)} \quad
          &\ket{x}\ket{y \xor \RO(x)}\ket{\RO(x) \xor \RO[G](x)}\, , \\
        \quad \apply{(5.)} \quad
          &\ket{x}\ket{y \xor \RO(x)}\ket{\RO(x) \xor \RO[G](x) \xor \RO(x)} \\
          &\quad = \ket{x}\ket{y \xor \RO(x)}\ket{\RO[G](x)}\, , \\
        \quad \apply{(6.)} \quad
          &\ket{x}\ket{y \xor \RO(x)}\ket{\RO[G](x) \xor \RO[G](x)} \\
          &\quad = \ket{x}\ket{y \xor \RO(x)}\ket{0^m}\, ,
        \intertext{
          at which point $\inputreg$ and $\outputreg$ are returned to $\ddv$.
          If $x \in \careset$, then we instead get,
        }
        \quad \apply{(4.)} \quad
          &\ket{x}\ket{y \xor \RO(x) \xor \RO(x) \xor \RO[G](x)}\ket{\RO(x) \xor \RO[G](x)} \\
          &\quad = \ket{x}\ket{y \xor \RO[G](x)}\ket{\RO(x) \xor \RO[G](x)}\, , \\
        \quad \apply{(5.)} \quad
          &\ket{x}\ket{y \xor \RO[G](x)}\ket{\RO(x) \xor \RO[G](x) \xor \RO(x)} \\
          &\quad = \ket{x}\ket{y \xor \RO(x)}\ket{\RO[G](x)}\, , \\
        \quad \apply{(6.)} \quad
          &\ket{x}\ket{y \xor \RO[G](x)}\ket{\RO[G](x) \xor \RO[G](x)} \\
          &\quad = \ket{x}\ket{y \xor \RO(x)}\ket{0^m}\, ,
      \end{align*}
      as desired.
      Crucially, we see that $\bdv$ ends up un-entangled with both
      $\distinguisher$ and
      the random oracles after returning $\inputreg$ and $\outputreg$ to
      $\distinguisher$. Thus, $\bdv$ acts only as an intermediary
      interface between $\distinguisher$ and the random oracles.
      \qed
    \end{proof}

    \begin{remark}
      It is also possible for $\bdv$ to employ the above strategy to
        adaptively reprogram $\altRO$ by changing the check in $\unitary_z$
        between queries, with the effect of altering the values in $\careset$.
        Of course, such behaviour may well be noticeable, and so requires a more
        careful analysis (such as that found in the proof of the adaptive
        reprogramming framework~\cite{PKC:PanZen24})\, .
    \end{remark}

    \begin{theorem}[Hands-Off O2H]\label{lemma:handsoff-o2h}
      Let $\RO_0: \bin^n \to \bin^m$ be a random oracle, and let distinguisher
        $\distinguisher$ be a quantum algorithm making at most $q$ queries to a
        random oracle. Let $\careset = \{ x^*_1, \ldots x^*_\ell \}$,
        and let $z$ be some value (or set of values) that depends arbitrarily
        on the points $\RO_0(x^*_1), \ldots, \RO_0(x^*_\ell)$. Let $\RO_1:
        \bin^n \to \bin^m$ be an independent random
        oracle. Then, $\distinguisher$'s distinguishing advantage against the
        \emph{hiding} game, in which it is given the input $z$ and oracle access
        to $\RO_\chalbit$, be 
      \[
        \advantage{\hide}{\RO_\chalbit, z}[\distinguisher] 
          \coloneqq 2 \cdot \prob{\chalbit \sample
          \distinguisher^{\RO_\chalbit}(z)} - 1\, .
      \]

      \noindent
      Now, define the extractor $\extractor$ to be such that it simulates $\RO_1$ for 
        $\distinguisher$, implemented as a compressed random oracle.
        As stated in \autoref{thm:perfect}, the simulation will be perfect, and so
        $\distinguisher$ will have the same distinguishing advantage in this case.

      After $\distinguisher$ halts with some output $\advbit$, let $\extractor$ 
        observe the database, obtaining a list $\Dlist$ of $q$ tuples $(x_i,
        \hat{y}_i)$. Then, define $\extractor$'s \emph{$1$-in-$q$ one-way}
        advantage as
      \[
        \advantage{q\hy\oneway}{z}[\extractor^\distinguisher] 
          \coloneqq \probsample{
            \{(x_i,\hat{y}_i)\}_{\forall i \in [q]} \sample
            \extractor^\distinguisher(z)
          }{
            \{x_i\}_{\forall i \in [q]} \cap \{x^*_i\}_{\forall i \in [\ell]} \neq \emptyset
          }\, .
      \]
      
      \noindent
      Then, for all distinguishers $\ddv$,
      \[
        \advantage{\hide}{\RO_\chalbit, z}[\distinguisher]
        \leq 2\sqrt{2 \cdot \advantage{q\hy\oneway}{}[\extractor^\distinguisher]}
        + \sqrt{2^{-m+3}}\, .
      \]
    \end{theorem}

    \noindent
    As you can see, we still end up with the same square root loss 
      so characteristic for O2H lemmas---but at least we are able to avoid the
      $q$ loss! This is essentially done by delegating the task of
      finding the right value out of $q$ options to the reduction
      applying the lemma.

    Whether this result is ``better'' than previous art will
      likely depend on the usecase, but for certain game hops 
      it will certainly be easier to apply, and may lead to tighter bounds when
      recognizing values that you care about is easy. On the other hand, there
      are clearly situations where the Resampling Lemma, for which the
      $q$ loss appears under the square root, will provide a tighter bound.

    What is most interesting to me, however, is seeing \emph{how} we
      seem doomed to run head-first into the inevitable square-root loss, 
      made explicit in the following proof.

    \begin{proof}
      \textbf{ERROR: \autoref{thm:imperfect} is not applicable.} 

      Assume without loss of generality that the adversary plays optimally, and
        further that they never perform any measurements (i.e.\ that
        $\distinguisher$ is
        unitary). Then, we may further assume that $\distinguisher$ will halt by outputting a
        single-qubit register $\hat{\register[B]}$, which $\distinguisher$ may have placed
        in one of three states:
        $\ket{0}$ resp.\ $\ket{1}$ if the adversary was able to distinguish (by
        querying $\RO_\chalbit$ on at least one point $x^*_i$), and
        the uniform superposition $\frac{\ket{0} + \ket{1}}{\sqrt 2}$ otherwise. 

      However, the state of $\hat{\register[B]}$ will of course be highly entangled with
        the state of $\distinguisher$ in general, and therefore also the database $\database$
        implementing $\RO_1$. What the above says, then, is that if we speciale
        to each basis state $\ket{\advbit}\ket{\distinguisher}\ket{\database}$ in the
        support of the combined superposition of the registers $\advbitregister,
        \distinguisher, \database$, we will have that in every branch where
        $\distinguisher$ queried
        $\RO_1$ on a value $x^*_i$, $\advbitregister$ is set to $\ket{1}$, as
        well as in half of the remaining branches, with $\advbitregister$ set
        placed in the $\ket{0}$ state in the other half.

      \noindent
      Let $\obsev$ be the event that $\distinguisher$ observed at least one
        $x^*_i$, and let $\alpha_\obsev$ be the amplitude of this happening,
        which may be isolated from the state using the projectors
        $\projector_\obsev$, $\projector_{\obsnev} = \idmatrix -
        \projector_\obsev$. Note that $\{\projector_\obsev,
        \projector_{\obsnev}\}$ do not correspond to a measurement
        we are actually able to perform as long we are treating $\distinguisher$
        as a black box; but as we will see, they \emph{do} correspond to an event
        that we are able to bound, in terms of a measurement we \emph{can} perform.

      Let $\ket{\advbitregister \distinguisher \database}_\obsev$ be state
        projected onto the subspace where $\obsev = \true$. Then,

      \begin{align*}
        \ket{\advbitregister \distinguisher \database} &= 
          \projector_\obsev \ket{\advbitregister \distinguisher \database}
          + (\idmatrix - \projector_\obsev) \ket{\advbitregister \distinguisher \database} \\
        &= \alpha_\obsev \ket{\advbitregister \distinguisher \database}_\obsev
          + \alpha_{\obsnev} \ket{\advbitregister \distinguisher
          \database}_{\obsnev} \\
        &= \alpha_\obsev \ket{\chalbit} \ket{\distinguisher \database}_\obsev
          + \alpha_{\obsnev} \left(
          \frac{\ket \chalbit + \ket{1-\chalbit}}{\sqrt 2}\right)
          \ket{\distinguisher \database}_{\obsnev} \\
        \intertext{
        where we inserted our assumption on $\distinguisher$'s strategy.
          We can write $\alpha_{\obsnev}$ in terms of its real and imaginary
          parts as $e^{i\theta_{\obsnev}} \sqrt{1 - \abs{\alpha_\obsev}^2}$, for
          some unknown phase angle $\theta_{\obsnev} \in [0, 2\pi)$:
        }
        \ket{\advbitregister \distinguisher \database} &= \alpha_\obsev
          \ket{\chalbit} \ket{\distinguisher \database}_\obsev
          + e^{i\theta_{\obsnev}}\sqrt{1 - \abs{\alpha_\obsev}^2} \left(
          \frac{\ket \chalbit + \ket{1-\chalbit}}{\sqrt 2}\right)
          \ket{\distinguisher \database}_{\obsnev} \\
        &= \alpha_\obsev \ket{\chalbit} \ket{\distinguisher \database}_\obsev
          + e^{i\theta_{\obsnev}} \sqrt{\frac{1 - \abs{\alpha_\obsev}^2}{2}}
          \ket{\chalbit} \ket{\distinguisher \database}_{\obsnev}
          + e^{i\theta_{\obsnev}} \sqrt{\frac{1 - \abs{\alpha_\obsev}^2}{2}} \ket{1-\chalbit}
          \ket{\distinguisher \database}_{\obsnev}
      \end{align*}

      \noindent
      The distinguisher wins
        iff.\ their output is observed to be $\chalbit$, and so we see from the
        above that their winning \emph{amplitude} is 

      \[
        \alpha_\distev \coloneqq \alpha_\obsev + 
        e^{i\theta_{\obsnev}} \sqrt{\frac{1 - \abs{\alpha_\obsev}^2}{2}}\, ,
      \]
      
      \noindent
      which, squared, gives us their distinguishing probability:
      \begin{align*}
        \prob{\distev} \coloneqq \abs{\alpha_\distev}^2 &= 
          \left( \alpha_\obsev + e^{i\theta_{\obsnev}} \sqrt{\frac{1 -
          \abs{\alpha_\obsev}^2}{2}} \right)^*
          \left( \alpha_\obsev + e^{i\theta_{\obsnev}} \sqrt{\frac{1 -
          \abs{\alpha_\obsev}^2}{2}} \right) \\
        &= \abs{\alpha_\obsev}^2 + \left( \alpha_\obsev^* e^{i\theta_{\obsnev}} 
          + \alpha_\obsev e^{-i\theta_{\obsnev}} \right) \sqrt{\frac{1 -
          \abs{\alpha_\obsev}^2}{2}} 
          + \frac{1 - \abs{\alpha_\obsev}^2}{2}\, .
      \end{align*}

      \paragraph{A failed attempt at tightness.}
        If we can get the middle terms to cancel out, then we will get a tight relation
          between the $\distinguisher$'s distinguishing advantage, and their
          probability of having observed an $x^*_i$, (which we can then extract):

        \begin{align*}
          \alpha_\obsev^* e^{i\theta_{\obsnev}} + \alpha_\obsev
          e^{-i\theta_{\obsnev}} &= 0\, .
          \intertext{
            Writing $\alpha_\obsev = \abs{\alpha_\obsev} e^{i \theta_\obsev}$ and
            factoring out the real part makes this
          }
          e^{- i \theta_\obsev} e^{i\theta_{\obsnev}} 
            + e^{i \theta_\obsev} e^{-i\theta_{\obsnev}} 
          %&= e^{i (\theta_{\obsnev} - \theta_\obsev)} 
          %  + e^{i (\theta_\obsev - \theta_{\obsnev})} 
          &= 0\, . \\
          \intertext{
            One natural thing to try is to multiply the state with a global
            phase $e^{i \theta}$, since such phases are unobservable anyway. This
            has the effect of shifting all phases by the angle $\theta \in [0,
            2\pi)$. What we want, then, is to find an angle $\theta$ such that
            the above holds.
          }
          e^{- i (\theta_\obsev + \theta)} e^{i (\theta_{\obsnev} + \theta)} 
            + e^{i (\theta_\obsev + \theta)} e^{-i (\theta_{\obsnev} + \theta)}
          &= e^{i (\theta - \theta + \theta_{\obsnev} - \theta_\obsev )} 
            + e^{i (\theta - \theta + \theta_\obsev - \theta_{\obsnev})} \\
          &= e^{i (\theta_{\obsnev} - \theta_\obsev )} 
            + e^{i (\theta_\obsev - \theta_{\obsnev})}\, .
          \intertext{
            A fool's hope. As we might have predicted, the global phase we
            introduced immediately cancelled out. We \emph{did} say it
            was unobservable, after all!
            This means that we're stuck with the following relation:
          }
          e^{i (\theta_{\obsnev} - \theta_\obsev )} 
            + e^{i (\theta_\obsev - \theta_{\obsnev})} 
          &= e^{i (\theta_{\obsnev} - \theta_\obsev )} 
            - e^{i (\theta_\obsev - \theta_{\obsnev} + \pi)} = 0\, , \\
          \intertext{
            which gives us that
          }
          \theta_{\obsnev} - \theta_\obsev = \theta_\obsev - \theta_{\obsnev} + \pi 
          %\quad &\Longrightarrow \quad 2 \theta_{\obsnev} = 2 \theta_\obsev + \pi \\
          \quad &\Longrightarrow \quad \theta_{\obsnev} - \theta_\obsev = \frac{\pi}{2}\, . \\
        \end{align*}
        In other words, we need the relative phase of $\alpha_{\obsev}$ and
          $\alpha_{\obsnev}$ to be \emph{exactly} $e^{i \nicefrac{\pi}{2}} = i$.
          This is not something we can guarantee will hold without
          making very precise assumptions on $\distinguisher$ (which we do not want
          to do, since the result is supposed to hold for all $\distinguisher$).

        What if it were? Well, then the middle terms would cancel out, and we would have

        \begin{align*}
          \prob{\distev} &= \abs{\alpha_\obsev}^2 + \frac{1 -
          \abs{\alpha_\obsev}^2}{2}
            = \frac{1}{2} \prob{\obsev} + \frac{1}{2}\, , \\
          \Longrightarrow \advantage{}{}[\distinguisher] &= 
            2 \cdot \left( \frac{1}{2} \prob{\obsev} + \frac{1}{2}\right) - 1 
            = \prob{\obsev}\, ,
        \end{align*}
        as desired. 

        %\hhnote{ Is this true? (Try setting $\theta_{\obsnev} = 0$.) \\
        % ``As an exercise, you may also verify that if we are given a guarantee
        %  such as $\alpha_{\obsnev}$ being a real number, then our strategy of
        %  rotating the global phase by an angle $\theta$ \emph{does} work.''
        %} A: Not true. In my earlier ``proof'', I had forgotten to take account of
        %   theta's effect on alpha_obsnev.

        In a way, it should not surprise us that there \emph{exists}
          distinguishers such that the relation is tight, since this is exactly
          what we get with classical distinguishers and lazy-sampling, of which
          quantum distinguishers and compressed oracles are a generalization.

        \hhnote{
          Can we show that tightness is achieved exactly when all queries are
          restricted to a single basis, so that they are equivalent to classical
          queries up to the choice of basis?
        }

      \paragraph{Giving in to the bound.}
        Let us return to what we know so far:

        \begin{align*}
          \prob{\distev} &= 
          \abs{\alpha_\obsev}^2 + \left( e^{i(\theta_{\obsnev} - \theta_\obsev)} 
            + e^{i( \theta_\obsev - \theta_{\obsnev}} \right) \abs{\alpha_\obsev} \sqrt{\frac{1 -
            \abs{\alpha_\obsev}^2}{2}} 
            + \frac{1 - \abs{\alpha_\obsev}^2}{2}\, .
          \intertext{
            Without being able to assume anything about $\theta_\obsev$ and $\theta_{\obsnev}$,
            we have no choice but to bound the middle term in terms of its
            highest possible value, which it reaches whenever $\theta_\obsev =
            \theta_{\obsnev}$:
          }
          \prob{\distev} &\leq
          \abs{\alpha_\obsev}^2 + 2 \cdot \abs{\alpha_\obsev} \sqrt{\frac{1 -
            \abs{\alpha_\obsev}^2}{2}} 
            + \frac{1 - \abs{\alpha_\obsev}^2}{2} \\
          &= \abs{\alpha_\obsev}^2 + \sqrt{
              2 \cdot \abs{\alpha_\obsev}^2 \left( 1 - \abs{\alpha_\obsev}^2 \right) 
            } + \frac{1 - \abs{\alpha_\obsev}^2}{2} \\
          &= \frac{1}{2} \prob{\obsev} + \sqrt{
              2 \cdot \prob{\obsev} \left( 1 - \prob{\obsev} \right) 
            } + \frac{1}{2}\, , \\
          \Longrightarrow \advantage{}{}[\distinguisher] 
            &= 2 \cdot \prob{\distev} - 1 \\
            &\leq \prob{\obsev} 
              + 2 \sqrt{ 2 \cdot \prob{\obsev} \left( 1 - \prob{\obsev} \right)
              }\, . 
        \end{align*}

        \noindent
        In order to give a clean bound in terms of the probability of
          observing one of the relevant values, we may at this point just as
          well throw away the smaller terms, and embrace the characteristic O2H
          square root.

        \[
          \advantage{}{}[\distinguisher] \leq 2 \sqrt{ 2 \cdot \prob{\obsev} }\, .
        \]

      \paragraph{Measuring the database.}
        Now we want to measure the database and extract a value $x^*_i$. Let
          $\extev$ be the event that the observed database list $\Dlist$
          contains at least one $x^*_i$, and define the binary projective measurement
          $\measurement_\extev$ corresponding to projectors
          $\{ \projector_\extev, \idmatrix - \projector_\extev \}$ which checks the
          input registers of the database, $\database_\inputreg$, and returns true
          iff.\ at least one element $(x^*_i, \hat{y})$ was observed. 

        Unlike the case of $\obsev$, this is a measurement we can
          actually perform without any assumptions on $\distinguisher$.
          It may for example be implemented as follows:
        \begin{enumerate}
          \item For a given subregister $\database_{\inputreg,j}$ of $\database$
            in some superposition over $x \in \bin^n \cup \emptystring$, flip the ancilla
            iff.\ the subregister houses the value $x^*_i$.
          \item Repeat (1) for all values $i \in [\ell]$, storing the result in a different
            ancilla each time.
          \item Compute the OR function between the ancillas 
            (which may be implemented using a series of $\Xgate$ and Toffoli
            gates), and store the result in a ``result'' ancilla, $\register[A]_j$.
          \item Repeat (1)--(3) for all $j \in [q]$.
          \item Compute the OR function between all result ancillas $\register[A]_j$, and
            store the result in one final ancilla $\register[A]$.
          \item Observe $\register[A]$ in the computational basis.
          \item Uncompute steps (1)--(5) by doing the operations in the reverse
            order (since the $\CXgate$, $\Xgate$, and Tofolli gates are all involutions).
        \end{enumerate}

      \noindent
      Let $\alpha_\extev$ be the amplitude resulting from applying the projector
        $\projector_\extev$ to the combined state $\ket{\advbitregister
        \distinguisher \database}$, or in other words, let the probability of
        $\measurement_\extev$ returning $1$ be $\prob{\extev} =
        \abs{\alpha_\extev}^2$. Then \autoref{thm:imperfect} tells us that 

      \[
        \sqrt{\prob{\obsev}} \leq \sqrt{\prob{\extev}} + \sqrt{2^{-m}}\, ,
      \]
      where we set the paramater $\ell$ in the theorem to $1$, since
        observing one $x^*_i$ is enough for $\obsev$ to be $\true$.

      One thing remains, namely to argue that performing the measurement
        $\measurement_\extev$ did not disturb the state, which could alter the result of
        observing the database and output register $\advbitregister$ in such a
        way that $\extev = \true$ would not be guaranteed to imply that observing
        the database gives one of the values we are looking for.

      To see this, note that our implementation of
        $\measurement_\extev$ only performed classical operations, as well as 
        one computational basis measurement. Theorem 2.2 in Nielsen \&
        Chuang~\cite[p.77]{NieChu10} implies that two projective measurements
        commute if and only if there exists
        an orthogonal basis such that the projectors defining the measurements are
        simultaneously diagonal. In our case, the computational basis is such a
        basis. This means that it does not matter if we perform the
        measurement $\measurement_\extev$ first and then collapse the output
        and database superpositions, or if we observe the output and database
        first and then check for a value $x^*_i$ using $\measurement_\extev$.

      Thus, we can be sure that if $\extev = \true$,
        observing the database will provide us with a list that includes at least
        one value $x^*_i$, which implies that $\prob{\extev} = \advantage{}{}[\extractor]$.
        Inserting the above inequality into our previous bound then gives us our
        final result:
        \begin{align*}
          \advantage{}{}[\distinguisher] &\leq 2 \sqrt{ 2 \cdot \prob{\obsev} } 
          \leq 2 \sqrt{2} \left( \sqrt{\prob{\extev}} + \sqrt{2^{-m}} \right) \\
          &= 2 \sqrt{2 \cdot \prob{\extev}} + 2\sqrt{2 \cdot 2^{-m}} 
          = 2 \sqrt{2 \cdot \advantage{}{}[\extractor]} + \sqrt{2^{-m+3}}\, .
        \end{align*}

      \qed 
      
      \hhnote{
        Q: Is there a simple-ish way to 
        get a slightly tighter bound from the fact that we are looking for one
        out of $\ell$ possible values, e.g.\ by apply the theorem several times, or adjusting its
        statement and proof? (Answer seems to be no.)
      }

    \end{proof}

    %\hhnote{12025-03-23: \\
    %  Still left to do: Implement a ``hop'' between $\prob{\distev}$ and $\prob{\extev}$ by
    %  applying \autoref{thm:imperfect}. \\
    %  Check if I need to specify the phase for $\alpha_{\extnev}$ to begin with.
    %  (This one could screw up the whole project.) % It did.
    %}

  \subsection{Joseph's Counterexample: An Uncomputing Adversary}
    \label{sec:uncompute}
    Consider the adversary $\adv$, given in \autoref{fig:uncomputing-adv}, that
      on input $(\target, \realRO(\target))$ and oracle access to $\idealRO$,
      trivially distinguishes $\realRO$ from $\idealRO$ by calling $\idealRO$
      on input $\target$ and comparing the output. Then, $\adv$ uncomputes its
      knowledge of the output $\idealRO(\target)$ by calling the oracle a
      second time, keeping only an ancilla
      qubit storing information on whether or not the values were equal.

    We want to know what the probability is of observing
      $\target$ in the database at the end of the protocol. We proceed by
      direct calculation of the combined quantum state.
      We will need the following simple, but useful, lemma.

    \begin{lemma}[Orthogonality Relation]\label{lemma:orthogonality}
      For all $n$-bit integers $x$, $y$, $z$,
      \[
        \sum_{x \in [2^n]} \frac{(-1)^{x \cdot (y \xor z)}}{2^n} = \delta_{y, z}\, ,
      \]
      where $\delta_{y, z}$ is the Kronecker delta.
    \end{lemma}

    \begin{proof}
      If $y = z$, then the exponent is zero independently of $x$, and the sum evaluates to
      $2^n$. For all other values of $y$ and $z$, there will be equally many
      values for $x$ that give a negative and a positive sign, respectively, so that the sum
      evaluates to $0$.
      \qed
    \end{proof}

    \noindent
    Let us now compute the state resulting from the operations given in
    \autoref{fig:uncomputing-adv}.
    Just before the first oracle query, we have the combined state
    \begin{align*}
      \ket{\bitreg}\ket{\inputreg}\ket{\outputreg}\ket{\database_0} &=
        \ket{0}_{\bitreg} \ket{\target}_{\inputreg} \ket{0^m}_{\outputreg}
        \ket{\emptystring, 0^m}_{\database_0}\, ,
      \intertext{
        where $\database_0$ is the database subregister that houses the first tuple
        in the list. As $\idealRO$ applies each unitary, as shown in
        $\autoref{fig:uncomputing-adv}$, this state is then transformed as
        follows.
      }
        %\ket{\bitreg}\ket{\inputreg}\ket{\outputreg}\ket{\database_0}
          &\apply{\decomp} 
          \ket{0}_{\bitreg} 
          \ket{\target}_{\inputreg} 
          \ket{0^m}_{\outputreg}
          \ket{\target, 0^m}_{\database_0} \\
        &\apply{\Hgate^{\otimes m}} 
          \ket{0}_{\bitreg} 
          \ket{\target}_{\inputreg} 
          \sum_{y \in [2^m]}
          \frac{1}{\sqrt{2^m}} 
          \ket{y}_{\outputreg}
          \ket{\target, 0^m}_{\database_0} \\
        &\apply{\entangle} 
          \ket{0}_{\bitreg} 
          \ket{\target}_{\inputreg} 
          \sum_{y \in [2^m]}
          \frac{1}{\sqrt{2^m}} 
          \ket{y}_{\outputreg}
          \ket{\target, y}_{\database_0} \\
        &\apply{\Hgate^{\otimes m}} 
          \ket{0}_{\bitreg} 
          \ket{\target}_{\inputreg} 
          \sum_{z \in [2^m]}
          \sum_{y \in [2^m]}
          \frac{(-1)^{z \cdot y}}{2^m} 
          \ket{z}_{\outputreg}
          \ket{\target, y}_{\database_0} \\
        &\apply{\decomp} 
          \ket{0}_{\bitreg} 
          \ket{\target}_{\inputreg} 
          \sum_{z \in [2^m]}
          \sum_{y \in [2^m] \backslash \{0^m\}}
          \frac{(-1)^{z \cdot y}}{2^m} 
          \ket{z}_{\outputreg}
          \ket{\target, y}_{\database_0} \\
        &\qquad + 
          \sum_{z \in [2^m]}
          \frac{1}{2^m}
          \ket{0}_{\bitreg} 
          \ket{\target}_{\inputreg} 
          \ket{z}_{\outputreg}
          \ket{\emptystring, 0^m}_{\database_0}\, .
        \intertext{
          At this point, $\idealRO$ returns $\inputreg$ and $\outputreg$ to
          $\adv$, who proceeds to do a conditional bitflip on $\bitreg$, conditioned on the
          state of register $\outputreg$ being $\ket{\targetout}$. Let us
          denote by $f(y)$ the predicate that is $=1$ when $y=\targetout$, and
          $=0$ otherwise. Then we have:
        }
        &\apply{\CXgate}
          \sum_{z \in [2^m]}
          \sum_{y \in [2^m] \backslash \{0^m\}}
          \frac{(-1)^{z \cdot y}}{2^m} 
          \ket{f(z)}_{\bitreg} 
          \ket{\target}_{\inputreg} 
          \ket{z}_{\outputreg}
          \ket{\target, y}_{\database_0} \\
        &\qquad + 
          \sum_{z \in [2^m]}
          \frac{1}{2^m}
          \ket{f(z)}_{\bitreg} 
          \ket{\target}_{\inputreg} 
          \ket{z}_{\outputreg}
          \ket{\emptystring, 0^m}_{\database_0}\, .
        \intertext{
          Now, $\adv$ makes its second query to $\idealRO$ with the same
          registers $\inputreg$, $\outputreg$, leading to the following
          transformations:
        }
        &\apply{\decomp}
          \sum_{z \in [2^m]}
          \sum_{y \in [2^m] \backslash \{0^m\}}
          \frac{(-1)^{z \cdot y}}{2^m} 
          \ket{f(z)}_{\bitreg} 
          \ket{\target}_{\inputreg} 
          \ket{z}_{\outputreg}
          \ket{\target, y}_{\database_0} \\
        &\qquad + 
          \sum_{z \in [2^m]}
          \frac{1}{2^m}
          \ket{f(z)}_{\bitreg} 
          \ket{\target}_{\inputreg} 
          \ket{z}_{\outputreg}
          \ket{\target, 0^m}_{\database_0} \\
        &\quad = 
          \frac{1}{2^m}
          \sum_{z,y}
          (-1)^{z \cdot y}
          \ket{f(z)}_{\bitreg} 
          \ket{\target}_{\inputreg} 
          \ket{z}_{\outputreg}
          \ket{\target, y}_{\database_0}\, ,
        \intertext{
          in which we let the range of the sums be implicit (they will always
          be over all of $\bin^m$ unless otherwise noted).
        }
        &\apply{\Hgate^{\otimes m}}
          \frac{1}{2^{\nicefrac{3m}{2}}}
          \sum_{w,z,y}
          (-1)^{y \cdot z \xor z \cdot w}
          \ket{f(z)}_{\bitreg} 
          \ket{\target}_{\inputreg} 
          \ket{w}_{\outputreg}
          \ket{\target, y}_{\database_0} \\
        &\apply{\entangle}
          \frac{1}{2^{\nicefrac{3m}{2}}}
          \sum_{w,z,y}
          (-1)^{y \cdot z \xor z \cdot w}
          \ket{f(z)}_{\bitreg} 
          \ket{\target}_{\inputreg} 
          \ket{w}_{\outputreg}
          \ket{\target, y \xor w}_{\database_0} \\
        &\apply{\Hgate^{\otimes m}}
          \frac{1}{2^{2m}}
          \sum_{v,w,z,y}
          (-1)^{y \cdot z \xor z \cdot w \xor w \cdot v}
          \ket{f(z)}_{\bitreg} 
          \ket{\target}_{\inputreg} 
          \ket{v}_{\outputreg}
          \ket{\target, y \xor w}_{\database_0} 
        \intertext{
          We would like to sum over $y$ and apply the orthogonality relation
          \autoref{lemma:orthogonality}. However, the state of the $\database_0$ register
          still depends on $y$. Towards remedying this, we make a change of
          variable from $w$ to $u \coloneqq y \xor w \Leftrightarrow w = y \xor u$.
        }
        &\quad =
          \frac{1}{2^{2m}}
          \sum_{v,(y\xor u),z,y}
          (-1)^{y \cdot z \xor z \cdot (y \xor u) \xor (y \xor u) \cdot v}
          \ket{f(z)}_{\bitreg} 
          \ket{\target}_{\inputreg} 
          \ket{v}_{\outputreg}
          \ket{\target, u}_{\database_0}\, .
        \intertext{
          Summing over all values of $y \xor u$, for fixed $y$, has the same
          effect as summing over all values of $u$, since the xor operation is
          a permutation. Therefore, we may write
        }
        &\quad =
          \frac{1}{2^{2m}}
          \sum_{v,z,y,u}
          (-1)^{y \cdot z \xor z \cdot y \xor z \cdot u \xor y \cdot v \xor u \cdot v}
          \ket{f(z)}_{\bitreg} 
          \ket{\target}_{\inputreg} 
          \ket{v}_{\outputreg}
          \ket{\target, u}_{\database_0} \\
        &\quad =
          \frac{1}{2^{2m}}
          \sum_{v,z,y,u}
          (-1)^{y \cdot (v \xor 0^m)}
          (-1)^{u \cdot (z \xor v)}
          \ket{f(z)}_{\bitreg} 
          \ket{\target}_{\inputreg} 
          \ket{v}_{\outputreg}
          \ket{\target, u}_{\database_0}\, ,
        \intertext{
          where we inserted a $0^m$ in the exponential to prepare for the next
          step. Now we sum over $y$, and apply \autoref{lemma:orthogonality}.
        }
        &\quad =
          \frac{1}{2^{m}}
          \sum_{v,z,u}
          \delta_{v,0^m}
          (-1)^{u \cdot (z \xor v)}
          \ket{f(z)}_{\bitreg} 
          \ket{\target}_{\inputreg} 
          \ket{v}_{\outputreg}
          \ket{\target, u}_{\database_0}\, .
        \intertext{
          Summing over $v$ gives
        }
        &\quad =
          \frac{1}{2^{m}}
          \sum_{z,u}
          (-1)^{z \cdot u}
          \ket{f(z)}_{\bitreg} 
          \ket{\target}_{\inputreg} 
          \ket{0^m}_{\outputreg}
          \ket{\target, u}_{\database_0}\, .
        \intertext{
          In order to analyze this expression, we split the expression into the
          two cases where $f(z) = 1$, which happens exactly when $z =
          \targetout$, and $f(z)=0$.
        }
        &\quad =
          \frac{1}{2^{m}}
          \sum_{u}
          (-1)^{\targetout \cdot u}
          \ket{1}_{\bitreg} 
          \ket{\target}_{\inputreg} 
          \ket{0^m}_{\outputreg}
          \ket{\target, u}_{\database_0} \\
        &\qquad +
          \frac{1}{2^{m}}
          \sum_{z \neq \targetout,u}
          (-1)^{z \cdot u}
          \ket{0}_{\bitreg} 
          \ket{\target}_{\inputreg} 
          \ket{0^m}_{\outputreg}
          \ket{\target, u}_{\database_0}\, .
        \intertext{
          In order to be able to use \autoref{lemma:orthogonality} again, we
          now add $0$ by adding and subtracting the $z = \targetout$ term, only
          with $\ket{\bitreg}= \ket 0$:
        }
        &\quad =
          \frac{1}{2^{m}}
          \sum_{u}
          (-1)^{\targetout \cdot u}
          \ket{1}_{\bitreg} 
          \ket{\target}_{\inputreg} 
          \ket{0^m}_{\outputreg}
          \ket{\target, u}_{\database_0} \\
        &\qquad +
          \frac{1}{2^{m}}
          \sum_{z,u}
          (-1)^{z \cdot (u \xor 0^m)}
          \ket{0}_{\bitreg} 
          \ket{\target}_{\inputreg} 
          \ket{0^m}_{\outputreg}
          \ket{\target, u}_{\database_0} \\
        &\qquad -
          \frac{1}{2^{m}}
          \sum_{u}
          (-1)^{\targetout \cdot u}
          \ket{0}_{\bitreg} 
          \ket{\target}_{\inputreg} 
          \ket{0^m}_{\outputreg}
          \ket{\target, u}_{\database_0}\, .
        \intertext{
          Summing over $z$ in the middle term, followed by $u$, then gives us
        }
        &\quad =
          \frac{1}{2^{m}}
          \sum_{u}
          (-1)^{\targetout \cdot u}
          \ket{1}_{\bitreg} 
          \ket{\target}_{\inputreg} 
          \ket{0^m}_{\outputreg}
          \ket{\target, u}_{\database_0} \\
        &\qquad +
          \sum_{u}
          \delta_{u,0^m}
          \ket{0}_{\bitreg} 
          \ket{\target}_{\inputreg} 
          \ket{0^m}_{\outputreg}
          \ket{\target, u}_{\database_0} \\
        &\qquad -
          \frac{1}{2^{m}}
          \sum_{u}
          (-1)^{\targetout \cdot u}
          \ket{0}_{\bitreg} 
          \ket{\target}_{\inputreg} 
          \ket{0^m}_{\outputreg}
          \ket{\target, u}_{\database_0} \\
        &\quad =
          \frac{1}{2^{m}}
          \sum_{u}
          (-1)^{\targetout \cdot u}
          \ket{1}_{\bitreg} 
          \ket{\target}_{\inputreg} 
          \ket{0^m}_{\outputreg}
          \ket{\target, u}_{\database_0} \\
        &\qquad +
          \ket{0}_{\bitreg} 
          \ket{\target}_{\inputreg} 
          \ket{0^m}_{\outputreg}
          \ket{\target, 0^m}_{\database_0} \\
        &\qquad -
          \frac{1}{2^{m}}
          \sum_{u}
          (-1)^{\targetout \cdot u}
          \ket{0}_{\bitreg} 
          \ket{\target}_{\inputreg} 
          \ket{0^m}_{\outputreg}
          \ket{\target, u}_{\database_0}\, .
        \intertext{
          There is a phase cancellation for the case $u=0^m$,
        }
        &\quad =
          \frac{1}{2^{m}}
          \sum_{u}
          (-1)^{\targetout \cdot u}
          \ket{1}_{\bitreg} 
          \ket{\target}_{\inputreg} 
          \ket{0^m}_{\outputreg}
          \ket{\target, u}_{\database_0} \\
        &\qquad +
          \left(
          1 - \frac{1}{2^m}
          \right)
          \ket{0}_{\bitreg} 
          \ket{\target}_{\inputreg} 
          \ket{0^m}_{\outputreg}
          \ket{\target, 0^m}_{\database_0} \\
        &\qquad -
          \frac{1}{2^{m}}
          \sum_{u \neq 0^m}
          (-1)^{\targetout \cdot u}
          \ket{0}_{\bitreg} 
          \ket{\target}_{\inputreg} 
          \ket{0^m}_{\outputreg}
          \ket{\target, u}_{\database_0}\, ,
        \intertext{
          and combining the two sums and factoring gives
          the final expression.
        }
        &\quad =
          \frac{1}{2^{m}}
          \ket{1}_{\bitreg} 
          \ket{\target}_{\inputreg} 
          \ket{0^m}_{\outputreg}
          \ket{\target, 0^m}_{\database_0} \\
        &\qquad +
          \left(
          1 - \frac{1}{2^m}
          \right)
          \ket{0}_{\bitreg} 
          \ket{\target}_{\inputreg} 
          \ket{0^m}_{\outputreg}
          \ket{\target, 0^m}_{\database_0} \\
        &\qquad +
          \frac{1}{2^{m}}
          \sum_{u \neq 0^m}
          (-1)^{\targetout \cdot u}
          \left(
            \ket{1}_{\bitreg} 
            - \ket{0}_{\bitreg} 
          \right)
          \ket{\target}_{\inputreg} 
          \ket{0^m}_{\outputreg}
          \ket{\target, u}_{\database_0} \\
        &\quad =
          \left(
            \left(
              1 - \frac{1}{2^m}
            \right) \ket{0}_{\bitreg} 
            + \frac{1}{2^{m}}
            \ket{1}_{\bitreg} 
          \right)
          \ket{\target}_{\inputreg} 
          \ket{0^m}_{\outputreg}
          \ket{\target, 0^m}_{\database_0} \\
        &\qquad +
          \frac{1}{2^{m}}
          \sum_{u \neq 0^m}
          (-1)^{\targetout \cdot u}
          \left(
            \ket{1}_{\bitreg} 
            - \ket{0}_{\bitreg} 
          \right)
          \ket{\target}_{\inputreg} 
          \ket{0^m}_{\outputreg}
          \ket{\target, u}_{\database_0}\, .
        \intertext{
          Finally, $\idealRO$ applies $\decomp$ one last time, 
        }
        &\quad =
          \left(
            \left(
              1 - \frac{1}{2^m}
            \right) \ket{0}_{\bitreg} 
            + \frac{1}{2^{m}}
            \ket{1}_{\bitreg} 
          \right)
          \ket{\target}_{\inputreg} 
          \ket{0^m}_{\outputreg}
          \ket{\emptystring, 0^m}_{\database_0} \\
        &\qquad +
          \frac{1}{2^{m}}
          \sum_{u \neq 0^m}
          (-1)^{\targetout \cdot u}
          \left(
            \ket{1}_{\bitreg} 
            - \ket{0}_{\bitreg} 
          \right)
          \ket{\target}_{\inputreg} 
          \ket{0^m}_{\outputreg}
          \ket{\target, u}_{\database_0}\, ,
    \end{align*}

    \noindent
    and returns $\inputreg$ and $\outputreg$ to $\adv$.

    The extractor can recover $\target$ iff
      $u \neq 0$. We get the probability that $u \neq 0$ by summing the
      amplitudes squared:
      \[
        \prob{\extev} = \frac{2^m - 1}{2^{2m}} 
        = \frac{1}{2^m} \left( 
          1 - \frac{1}{2^{m}} 
        \right)
        < \frac{1}{2^m}\, .
      \] 

    \noindent
    The distinguisher, meanwhile, has probability 
    \begin{align*}
      \condprob{\distev}{u = 0} &= \left( 
        1 - \frac{1}{2^m} \right)^2
        = 1 - \frac{2}{2^m} + \frac{1}{2^{2m}} 
        > 1 - \frac{1}{2^{m-1}}
      \intertext{
        of distinguishing when $u\neq 0$, and probability $\half$ otherwise,
        leading to the total distinguishing probability
      }
      \prob{\distev} &= 
        \prob{u = 0}
        \condprob{\distev}{u = 0}
        + \prob{u \neq 0}
        \condprob{\distev}{u \neq 0} \\
      &> 
        \left( 1 - \frac{1}{2^m} \right) 
        \left( 1 - \frac{1}{2^{m-1}} \right)
        + \frac{1}{2^m} \left( 1 - \frac{1}{2^{m}} \right) 
        \frac{1}{2}
        > 1 - \frac{1}{2^m}\, .
    \end{align*}

    \paragraph{Conclusion:} An adversary really does have the ability to thwart the efforts
      of any extractor simply by uncomputing its oracle queries.
      Therefore, it seems aborting and measuring a random query really is
      required in order to have any hope of extracting.

    \paragraph{Next question:} Can we use this to prove the optimality of O2H?
      Could phrase counterexample in terms of a KEM scheme just hashes the
      keys, by starting from $\owqpca$ and letting $\adv$ be the one that
      performs measurements of the type $\measurement_\obsev$ on its own
      collection of queries, and uncomputing after each query.

  \subsection{Connecting with the Real World}
    \hhnote{
      In which I introduce the more fine-grained analysis of counting depth and
      breadth separately, as well as quantum vs.\ classical oracle calls. Thus,
      if your loss is only in the depth of quantum calls to the oracle, one can
      argue that this is much smaller than the $2^{80}$ loss that one would
      arguably get from just a linear $q_{\RO}$ loss. (Indeed, in the current NISQ-era,
      it will be \emph{very} small! Like, it is currently on the order of
      \emph{maybe} $2^3$, for sufficiently small input and output domains.)
    }
  \fi
